

\section{Thirtieth Sunday after Trinity: Brotherly Love}


\begin{quote}

John 13:34–35. A new commandment I give to you, that you love one another; that, as I have loved you, you also love one another. By this shall all know that you are my disciples, if you have love among yourselves.

\end{quote}

\bigskip

Jesus Christ’s new commandment is given to Jesus Christ’s disciples and applies to them, and to them alone. For the new commandment has its root and ground in these words of our text: "As I have loved you." Jesus’ love toward us is the fountain of our love toward the brethren. It is also the measure by which we are to measure, and according to which we are to test, our brotherly love.

Therefore this is the first probing question which our text awakens in our souls: Have we in truth experienced and known Jesus Christ’s love toward us? For without a living experience of Jesus’ love, it is of no use to try to love the brethren. It would only become a new bondage under the law and a new compulsion; but compelled love is no true and genuine love. It would only become a yoke which no one could bear.

How, then, can anyone come to the living experience of Jesus’ love? You have surely heard that often, and you ought to know it quite well. Even if misleading voices arise and say, Jesus loves whom he will, and you can do nothing but wait until his love seizes you; yet you know better, Soul, than to listen to such seductive words. You surely know that Jesus stands at the door of the heart and knocks and wills that you should open to him. You surely know that Jesus’ love has been and is and will continue to be toward you, so long as it is the time of grace and the day of salvation. It is not because Jesus does not love you, does not seek you, that you do not experience his love, but because you keep the door shut fast when he knocks, keep the heart closed when he seeks you.

How, then, is the door of the heart opened to the Savior? That too you surely know, if only you were willing to do what you know so well. You keep the heart closed and the door shut fast because you know that your heart is unclean, because you understand that sin must go when Jesus comes in. Then there is no other way by which you can open the heart’s door to the seeking, knocking Savior than by the confession of your sins. Do as David says: "I acknowledged my transgressions and did not hide the iniquity of my sin," and the door is opened. Fall down as Peter did at Christ’s feet and say: "Lord, I am a sinful man," and the heart is open. Strike your breast with the publican and pray: "God, be merciful to me, a sinner," and there is room in the heart for Jesus’ saving grace.

That is the way to come to the living experience of Jesus’ love. Without confession of sins there is no forgiveness of sins; but in the forgiveness of sins there is the living experience of Jesus’ love. He died for you because you were a sinner, fallen under death and judgment; as Paul says: "The death that he died, he died once for sin;" and as we read in the Epistle to the Hebrews: "Christ was once offered to bear away the sins of many." His death is his love, and his love is his death; but his death is for your sins, Soul, and therefore you have no whole and full experience of his love until you confess your sin to him. Then he enters into the unclean house of the heart and speaks his mighty word: "Son, Daughter, take heart; your sins are forgiven you."

Then you experience that he loves you, loves you so highly that he forgives you your sin and cleanses you from all unrighteousness; loves you so highly that he gives you life from his life and Spirit from his Spirit. Then you learn to love him who first loved you.

If this is your experience, then you know the love of Christ, which surpasses knowledge. Then the Spirit also teaches you to love the brethren. Or how should anyone be full of Christ’s love and Christ’s life and not love those who belong to Christ?

But the second probing question which our text awakens in a Christian heart is this: Do I love the brethren just as Christ has loved me? As highly, as deeply, as inwardly, and as self-sacrificingly? How does it stand, Brother, when your love is to be measured by this measure? Do you live for the brethren, just as he lived for them? Are you willing to die for the brethren, as he laid down his life for his friends? Do you live in such a way for the brethren that, when death comes, you can say in truth: I have spent my strength in the labor for God’s church; I have loved it unto death.

If that which is called Christianity among us is to be measured by this measure, how much of it is there then that remains? We grumble and complain when God’s church requires a very little sacrifice, a very small outlay. We are ready to be unkind toward those who go about gathering the money that is needed for the work of God’s kingdom. How does this accord with brotherly love, with love for Jesus and for his church?

Brethren! this ought not so to be. Our hearts and our hands ought to be opened for God’s church. Jesus has bought it with his blood; how should we not love it in heart and deed? When God’s church has need of our money, our labor, our self-denial, our sacrifice, we ought to hasten to bring what we are able. And when the messengers of God’s church come to us and say: Give me what is needed for the church’s work, then we ought to thank them with glad hearts because they do what we were duty-bound to do.

Let this become better hereafter. Think on Christ’s love, how heavy and narrow and bloody was the way that he went, in order that God’s church might be built among us, and love the church and the brotherhood as he loved. You must take this matter in earnest, if your Christian life is not to wither and die. What is the church, if it is not full of love? What is brotherhood, if there is no brotherly love?

And Jesus’ commandments are not burdensome; his yoke is good, and his burden is light. If a human heart is enlarged by Jesus’ love, then it is not heavy and painful to love God’s church and Jesus Christ’s brethren. It is the death of the flesh, that is true; but for the Spirit it is life and joy. Blessed is the heart that loves, and the life that is offered up in love.

But it is not only for our own sake, for the brotherhood’s sake, that we are to love the brethren. Jesus wills that we should love one another, so that all may know that we are his disciples. It belongs to the calling of Christians in the world to bear witness and confess and reveal God’s love to the world. And this takes place through mutual love. "See how they love one another;" that is a revelation of God’s love which we owe to all. Our fellowship with God the Father is to be in secret, that the hidden man of the heart may thrive. The roots are hidden in the bosom of the earth, but trunk and crown grow up gloriously in the clear air. So too must the life-roots of a child of God be hidden in God. But brotherly love must shoot forth in rich growth before the eyes of all, that they may see and wonder at that power of God which shows itself mightily in his children. If God’s children loved one another more, the coldness and darkness of the world would have to give way more before the power of love. The world says to Christians: Do miracles, and we will believe. That is the same thing the Jews said to Jesus: Come down from the cross, and we will believe. But this is Jesus Christ’s miracle, that he did not come down from the cross, but suffered death in love toward the world. Thereby he became victor over the world and death and the devil. That is also the mode of operation of God’s church. Remain faithful in love, and your whole life is a miracle in the midst of the world’s hatred and wickedness and self-love. Desire not yourself to do any other miracle than this. And do not think that the power of unbelief is overcome by any other miracle than this. Lay all your diligence upon love, and the light of love shall shine much farther than you can expect or understand.

Brethren! let us love, not in word, neither with the tongue alone, but in deed and truth!
